% Pursuing Stacks (cotes 134-2 to 134-8): archive cote 134-5, site page 37 (PDF page 38), typescript page 224.
% typescript. Draft compared: ps4.tex:1791-1864.
% AI-transcribed by Claude Opus 5.5 (claude-opus-5-5, adaptive thinking, effort high), Message Batches API, batch ps2 (msgbatch_01NVwr1Yuo8qJf5rykeGBXPK), 2026-11-12; not yet checked by eye.
\scanpage{134-5/37}
\mspage{224}
to $\overline{i_{B}}$ , using the \xout{fxx} canonical functor
\[ j_B = i_B^{*} : \ (Cat) \longrightarrow B^{\wedge} \ , \qquad X \longrightarrow B_{/X} \ \xout{\unsure{xxHom}} \ . \]
Thus we get a description of an equivalence
(\over{**}{\unsure{1'}}) \[ \overline{j_B}\ \overline{i_A} \ : \ Hot_A \xrightarrow{\hand{\approx}} Hot_B \ , \]
whose quasi-inverse of course is just the similar $\overline{j_A}\ \overline{i_B}$ . We can vary
a little this description, admittedly cumbersome in practice, by repla-%
cing $\overline{\over{\unsure{i}}{j}_{B}}$ by the isomorphic functor $\overline{i^{*}}$, where $i : B \to (Cat)$ is
any test functor from B to (Cat) (while we have to keep however $\overline{i_A}$ as it
is , without the possibility of replacing $i_A$ by a more amenable test
functor).

Another way for comparing $Hot_A$ and $Hot_B$ arises, as we saw yester%
day, whenever we have an \uline{aspheric} functor
\hand{(2)} \[ i : A \longrightarrow B \ , \]
by just taking
\hand{(2')} \[ \overline{i^{*}} \ : \ Hot_B \xrightarrow{\hand{\approx}} Hot_A \ . \]
This of course is about the \xout{xxxx} simplest way imaginable, all the more
as the functor $i^{*} : B^{\wedge} \to A^{\wedge}$ commutes to arbitrary direct and inverse
limits, just perfect for comparing constructions in $B^{\wedge}$and constructions
in $A^{\wedge}$ - whereas the functor $j_B i_A : A^{\wedge} \longrightarrow B^{\wedge}$ , giving rise to (*) above,
commutes just to sums and fibered products, not to amalgamated sums nor
to products, sad enough! We could add here that if we got \uline{two}
aspheric functors from A to B, namely i plus
\[ i' : \ A \longrightarrow B \ , \]
then (as immediately checked) any map between these functors
\[ u : \ i \longrightarrow i' \]
gives rise to an \xout{equivalence}\inserted[typed]{isomorphism} between the corresponding equivalences
\diagram{\begin{tikzcd}[column sep=large]
Hot_B \ar[r, bend left, "\overline{i^{*}}", ""{name=U, below}] \ar[r, bend right, "\overline{i'^{*}}"', ""{name=D, above}] & Hot_A \ar[from=D, to=U, Rightarrow, "\hand{\overline{u}}"']
\end{tikzcd}} ,

which is nothing but the canonical\xout{x}isomorphism referred to yesterday
\margin{(*)} (both functors being canonically isomorphic to (\over{**}{\unsure{1'}})), and yields the
most evident way for \xout{x}``computing'' the latter. Thus it turns out that
the isomorphism $\overline{u}$ does not depend upon the choice of u. If we want
to ignore this fact and look at the situation sternly, in a wholly compu-%
tational spirit, we could present things by stating that we get a
contravariant functor, from aspheric maps $i: A \to B$ to equivalences
$Hot_B \to Hot_A$ :
\margin{37}
